% Environmental exposure section.
% Numbers are sourced from:
%   notebooks/env_association_study.ipynb
%   notebooks/env_forecasting_value.ipynb
%   outputs/env_association_summary.md
%   outputs/env_forecasting_summary.md
%   outputs/env_temperature_heterogeneity_summary.md
%   outputs/env_gate_survivors.csv
%   outputs/env_residual_head_results.csv
% No causal claims are made.

\section{Environmental Exposure as an External Modality for Overnight Glucose Stability}
\label{sec:environmental_exposure}

\subsection{Motivation}
\label{subsec:env_motivation}

The previous sections evaluate information streams that are internal to the
forecasting task: CGM history, wearable signals, sleep context, static
participant descriptors, and derived event states. This section asks whether a
more external modality, the bedroom environment, contains additional
information about overnight glucose stability. The analysis was intentionally
staged. First, environmental exposure was tested as a descriptive candidate
signal associated with overnight glycemic outcomes. Second, the only surviving
association was tested for forecasting value using a frozen SSM-CGM residual
head. This separation matters: an association with a glucose summary does not
by itself justify adding a new input channel to the forecasting model.

\subsection{Construction of Bedroom-Valid Exposure Nights}
\label{subsec:env_exposure_nights}

Environmental sensors were fixed in the home, so a raw sensor stream is not
automatically a personal exposure stream. The primary analysis therefore kept
only high-confidence bedroom or bedside placements and excluded living room,
kitchen, office, hallway, and unknown placements. Sleep windows were derived
from wearable sleep-stage signals, and a valid night required overlapping
sleep, CGM, and environmental coverage. This filtering was required because a
fixed home sensor outside the bedroom may measure the dwelling environment
rather than the participant's sleep exposure.

\begin{table}[H]
\centering
\caption{Cohort construction for bedroom environmental exposure analysis.}
\begin{tabular}{l r}
\toprule
Filtering step & Count \\
\midrule
Selected AI-READI cohort & 1591 participants \\
High-confidence bedroom placement & 574 participants \\
Exposure-valid participants & 566 participants \\
Exposure-valid nights & 4419 nights \\
Median nights per participant & 9 nights \\
Mean nights per participant & 7.81 nights \\
\bottomrule
\end{tabular}
\label{tab:env_cohort_construction}
\end{table}

\begin{table}[H]
\centering
\caption{Environmental exposure and overnight glucose outcome groups.}
\begin{tabular}{p{0.28\linewidth} p{0.62\linewidth}}
\toprule
Group & Variables \\
\midrule
Bedroom environment &
\texttt{temp\_mean} (mean temperature), \texttt{temp\_sd} (temperature
variability), \texttt{hum\_mean} (mean relative humidity),
\texttt{pm25\_auc} (integrated particulate exposure),
\texttt{voc\_spike\_burden}, \texttt{nox\_spike\_burden}, and
\texttt{night\_light\_exposure}. \\
Overnight glucose &
Mean glucose, glucose standard deviation, coefficient of variation, glucose
range, time below 70 mg/dL, time above 180 mg/dL, wake-to-morning rise, and
dawn-associated rise. \\
\bottomrule
\end{tabular}
\label{tab:env_exposure_outcome_groups}
\end{table}

\subsection{Within-Participant Association Model}
\label{subsec:env_association_model}

Each exposure-outcome pair was fit with a within-between mixed-effects model:
\begin{equation}
Y_{i,n}
=
\beta_W (x_{i,n}-\bar{x}_i)
+
\beta_B \bar{x}_i
+
u_i
+
\epsilon_{i,n},
\label{eq:env_within_between}
\end{equation}
where $Y_{i,n}$ is the glycemic outcome for participant $i$ on night $n$,
$x_{i,n}$ is the environmental exposure on that night, $\bar{x}_i$ is the
participant-level mean exposure, $x_{i,n}-\bar{x}_i$ is the within-participant
exposure deviation, and $u_i$ is a participant random intercept. The
coefficient $\beta_W$ asks whether the same participant has a different
overnight glucose outcome on nights when their bedroom environment differs
from their own usual environment. This target is more appropriate than a simple
correlation because it reduces confounding by stable participant differences.

Because the grid contained many exposure-outcome pairs, raw $p$-values were not
used as the decision rule. Benjamini-Hochberg FDR correction was applied across
the primary exposure-outcome grid, and participant-clustered bootstrap
confidence intervals were used because one participant can contribute multiple
nights. A result was considered robust only if the FDR-corrected $q$-value was
below 0.05, the participant-bootstrap 95\% confidence interval excluded zero,
and temporal negative controls did not reproduce the association. Shuffled-night
and time-shift controls are the strongest controls in this setting. Non-bedroom
and daytime analyses were treated only as specificity checks because personal
proximity to the sensor is unknown outside bedroom sleep periods.

\subsection{Primary Environmental Association Results}
\label{subsec:env_primary_results}

Across 56 primary exposure-outcome pairs, the only pair that cleared the gate
was mean bedroom humidity against overnight glucose range:
\begin{equation}
\texttt{hum\_mean} \rightarrow \texttt{overnight\_glucose\_range}.
\end{equation}
The within-participant coefficient was $-0.361$ mg/dL per 1 percentage-point
higher humidity than participant usual (FDR $q=0.0466$; participant-bootstrap
95\% CI $[-0.600,-0.118]$; 566 participants; 4419 nights). In words, on nights
when a participant's bedroom humidity was higher than their own average, their
overnight glucose range tended to be narrower. The effect size is modest: a
10 percentage-point humidity difference corresponds to approximately
3.6 mg/dL lower overnight glucose range under the linear model. This is a
glucose stability association, not evidence for a mean-glucose effect and not a
causal estimate.

\begin{figure}[H]
  \centering
  \includegraphics[width=0.92\linewidth]{generated/env_humidity_association.png}
  \caption{Primary within-participant association between bedroom humidity and
  overnight glucose range. Higher-than-usual bedroom humidity was associated
  with a narrower overnight glucose range, and the participant-bootstrap
  confidence interval excluded zero. Shuffled-night and +24 h shifted exposure
  assignments attenuated the association and had confidence intervals crossing
  zero.}
  \label{fig:env_humidity_association}
\end{figure}

\subsection{Robustness, Nonlinear, and Heterogeneity Analyses}
\label{subsec:env_robustness}

The humidity finding passed several follow-up checks. A humidity bounds audit
found no invalid values outside the plausible 0--100\% range. After trimming
the most extreme 1\% of within-person humidity deviations, the coefficient
strengthened from about $-0.361$ to $-0.441$, with the bootstrap confidence
interval still excluding zero. The real same-night coefficient was not
reproduced by shuffled-night exposure assignment and was strongly attenuated
under the +24 h shifted exposure control. The 12 h shift is not emphasized
because its implementation reproduced the original alignment. Non-bedroom and
daytime comparisons were used only as supportive specificity checks, not true
negative controls.

Secondary nonlinear models tested whether temperature or humidity had a hidden
U-shaped or absolute-deviation relationship:
\begin{equation}
Y \sim (x-\bar{x}) + (x-\bar{x})^2 + \bar{x} + u_i + \epsilon,
\end{equation}
and
\begin{equation}
Y \sim |x-\bar{x}| + \bar{x} + u_i + \epsilon.
\end{equation}
Across 48 nonlinear tests, zero were FDR-significant and robust. No quadratic
or absolute-deviation effect was found for temperature or humidity, supporting
the simpler linear within-person model for the main humidity result.

Temperature was then evaluated as a possible heterogeneous exposure rather than
a single cohort-wide effect. The subgroup interaction grid contained 204
temperature tests; three model-based FDR-significant signals appeared for
\texttt{temp\_sd} against overnight time above 180 mg/dL, involving insulin
use, top BMI quartile, and top sleep-fragmentation quartile. All participant
bootstrap intervals included zero, so these were labeled model-based only and
not participant-robust. Participant-level temperature slopes were similarly
exploratory: among 566 participants, 102 were classified as positive
temperature-sensitive candidates, 102 as negative candidates, 303 as not
sensitive, and 59 had too few nights. Held-out early/late validation showed
63\% sign agreement and near-zero slope correlation across 471 participants.
Temperature therefore remains a hypothesis-generating heterogeneous exposure,
not a validated personalization signal.

\begin{figure}[H]
  \centering
  \includegraphics[width=0.72\linewidth]{generated/env_temp_sensitivity_shrunk_slope_hist.png}
  \caption{Exploratory participant-level temperature sensitivity estimates.
  Candidate positive and negative sensitivity groups were observed, but
  early-late night validation showed weak reproducibility, so these slopes were
  not treated as stable personalized phenotypes.}
  \label{fig:env_temperature_sensitivity}
\end{figure}

The notebook also profiled the candidate groups against clinical and medication
context (Figs.~\ref{fig:env_temperature_clinical_boxplot}--\ref{fig:env_temperature_medication_studygroup}). These plots are descriptive checks on whether the slope-derived groups were clinically interpretable. Positive temperature-sensitive candidates tended to have higher baseline glucose variability and somewhat higher HbA1c and BMI, while the negative and non-sensitive groups were closer together. Medication and study-group composition also differed across candidate groups, especially for insulin use and diabetes-status strata. This pattern makes the candidate labels biologically plausible enough to report, but it also reinforces the main caveat: the labels partly track baseline glycemic instability and cohort composition, and the held-out slope validation was weak. They therefore support hypothesis generation, not a deployed personalization rule.

\begin{figure}[H]
  \centering
  \includegraphics[width=0.95\linewidth]{generated/env_temp_sensitivity_clinical_boxplot.png}
  \caption{Clinical variables by exploratory temperature-sensitive candidate group for \texttt{temp\_sd} versus overnight time above 180 mg/dL. The candidate groups show descriptive differences in HbA1c, BMI, baseline glucose variability, and awakening count, but these differences are interpreted as characterization of an unstable slope label rather than validation of a participant phenotype.}
  \label{fig:env_temperature_clinical_boxplot}
\end{figure}

\begin{figure}[H]
  \centering
  \includegraphics[width=0.95\linewidth]{generated/env_temp_sensitivity_insulin_studygroup_bar.png}
  \caption{Medication use and study-group composition by exploratory temperature-sensitive candidate group. Insulin use and diabetes-status strata differ across candidate groups, suggesting that apparent temperature sensitivity may partly reflect clinical subgroup structure and baseline glycemic instability.}
  \label{fig:env_temperature_medication_studygroup}
\end{figure}

After the humidity association was identified, sleep architecture was added to
test whether the association was simply explained by measured sleep summaries.
This was an adjustment and attenuation analysis, not a formal causal mediation
analysis:
\begin{equation}
Y_{i,n}
=
\beta_H H^{\mathrm{within}}_{i,n}
+
\beta_B \bar{H}_i
+
\theta^\top S_{i,n}
+
u_i
+
\epsilon_{i,n},
\label{eq:env_sleep_adjusted}
\end{equation}
where $S_{i,n}$ contains sleep summaries such as sleep duration and
sleep-stage composition. The humidity coefficient did not shrink toward zero:
it moved from $-0.361$ before sleep adjustment to $-0.391$ after adjustment,
and no significant exposure-by-sleep interaction was detected. Thus the
association was not explained away by measured sleep architecture, while still
remaining a descriptive association rather than proof of independent sleep
physiology.

\subsection{From Association to Forecasting Value}
\label{subsec:env_forecasting_value}

A statistically robust association is not enough to justify increasing the
SSM-CGM input space. The forecasting analysis therefore froze the base SSM-CGM
model and tested whether humidity explained residual forecast error. Let the
frozen model produce
\begin{equation}
\hat{y}^{\mathrm{base}}_{i,t+h}.
\end{equation}
The residual is
\begin{equation}
r_{i,t+h} = y_{i,t+h} - \hat{y}^{\mathrm{base}}_{i,t+h}.
\end{equation}
A lightweight residual head was trained as
\begin{equation}
\Delta \hat{y}_{i,t+h}
=
g_{\phi}
\left(
\hat{y}^{\mathrm{base}}_{i,t+h},
H_{i,n},
S_{i,n},
c_{i,t}
\right),
\end{equation}
and the corrected forecast was
\begin{equation}
\hat{y}^{\mathrm{corrected}}_{i,t+h}
=
\hat{y}^{\mathrm{base}}_{i,t+h}
+
\Delta \hat{y}_{i,t+h}.
\end{equation}
Here $H_{i,n}$ is the night-level humidity feature, $S_{i,n}$ denotes sleep or
window context, $c_{i,t}$ denotes available base forecast context, and
$g_{\phi}$ is a lightweight ridge or histogram-gradient residual model. This
design isolates the incremental value of humidity: any gain must come from
explaining the frozen base model's remaining error, rather than from full
retraining, optimization noise, or representation changes.

\begin{table}[H]
\centering
\caption{Forecasting-value comparison for humidity residual correction.}
\begin{tabular}{l l}
\toprule
Model & Purpose \\
\midrule
Frozen SSM-CGM base & Reference forecast \\
Base + sleep & Tests sleep-only residual value \\
Base + humidity & Tests environmental residual value \\
Base + sleep + humidity & Tests humidity beyond sleep context \\
Base + shuffled humidity & Negative control for spurious environmental gain \\
\bottomrule
\end{tabular}
\label{tab:env_forecast_models}
\end{table}

The 30-minute early-wake HGBT residual head was the clearest real-humidity
case: \texttt{base\_plus\_env} improved MAE by only 0.046 mg/dL on held-out
participants, with a 95\% confidence interval crossing zero
($[-0.061,0.155]$ mg/dL). This is more than six times smaller than the
pre-defined 0.3 mg/dL threshold for a meaningful forecasting gain. It also did
not beat the corresponding shuffled-night humidity control on the same
30-minute test endpoint; the shuffled control showed a numerically larger
0.066 mg/dL gain with a confidence interval also crossing zero. The notebook's
final decision file selected \texttt{standalone\_attribution\_only}: no tested
horizon cleared the full candidate gate, and humidity should not be added as an
SSM-CGM input channel at this time.

\begin{figure}[H]
  \centering
  \includegraphics[width=0.92\linewidth]{generated/env_forecasting_value.png}
  \caption{Forecasting-value test for bedroom humidity. Although humidity was
  associated with overnight glucose range in the association analysis, a
  residual head using humidity produced only a negligible MAE gain and did not
  outperform the shuffled-humidity control on the highlighted 30-minute
  early-wake endpoint.}
  \label{fig:env_forecasting_value}
\end{figure}

Several implementation limits are relevant to the forecasting-value test.
Sections 4--8 of the forecasting notebook used the
\texttt{experiment\_c\_split\_adapt6h\_seed42} split because a validated dense
prediction cache was already available, while the association study used the
\texttt{adapt48h\_seed42} split. The decisive comparison between real humidity
and shuffled humidity remains internal to the forecasting split. No fresh full
inference was run; the existing dense forecast cache was filtered and reshaped.
Only short-horizon residual correction was evaluated, and the 30-minute
early-wake endpoint was the clearest candidate display. Given the tiny observed
gain and the failure against shuffled humidity, these deviations are unlikely
to change the decision, but they constrain the claim.

\subsection{Interpretation and Decision}
\label{subsec:env_decision}

Together, these analyses separate three levels of evidence. First, bedroom
humidity showed a statistically robust within-participant association with
overnight glucose range. Second, this association was not explained away by
measured sleep architecture and was robust to key temporal controls. Third,
when tested as an input to a residual forecasting head on top of frozen SSM-CGM
predictions, humidity produced only a negligible MAE improvement and failed
against its shuffled-night control. Therefore, humidity is retained as an
interpretable, attribution-only environmental association, but it is not added
as an SSM-CGM input channel.
